\section{Logistic core: crossing, precision, and parameter sensitivity}\label{si:core}

\subsection{Equilibrium and linear map}

The logistic core uses \(S(N)=rN(1-N/K)\) and the zero-deficit equilibrium \(Z^*=\delta=\log(2)/10\). Solving the interior root of the effort law yields \(E^*=2.089623397737668\); the stock equilibrium is \(N^*=K(1-qE^*/r)=89.55188301131166\). The equilibrium fishing mortality is \(F^*=qE^*=0.002089623397737668\,\mathrm{yr}^{-1}\). Local effort gains are \((C_E,C_Z)=(-0.0595180098426,1.7850187149850)\) for the mobilising law and \((-0.850336140482,-1.661701896635)\) for the protective comparator.

For held effort, the plant's linearisation is the matrix \(A_{\rm hold}\) retained in \texttt{\detokenize{work/results.json}} and regenerated by \texttt{\detokenize{work/p07_core.py}}. The review matrix composes the held plant flow with either (i) the exponential frozen-assessment controller-flow update, whose local linearization is exact for the controller ODE with the endpoint assessment fixed, or (ii) the separately defined forward-Euler update. Eigenvalues are calculated for the product in the registered flow-then-update ordering. The full state dimension is three; no continuous-delay approximation is used in this matrix calculation.

\subsection{Baseline operator record}

\begin{table}[htbp]
\centering\small
\caption{Baseline radii at annual review and all computed crossings on the interval scan. Values are from the retained full-precision core output.}
\label{tab:core-crossings}
\begin{tabularx}{\linewidth}{@{}l r X@{}}
\toprule Map & \(\rho(1\,\mathrm{yr})\) & Crossing record / grid interpretation \\
\midrule
Mobilising, exponential frozen assessment & 1.0003506565 & One complex-pair crossing at 6.5012949481 yr, unstable \(\to\) stable. \\
Mobilising, Euler command step & 1.0005452356 & Complex-pair crossing at 47.5359822758 yr, unstable \(\to\) stable; real \(-1\) crossing at 79.1426878068 yr, stable \(\to\) unstable. \\
Protective, exponential frozen assessment & 0.9838095183 & No crossing on \([0.2,200]\) yr; maximum sampled \(\rho=0.996740\). \\
Protective, Euler command step & 0.9837962340 & Real \(-1\) crossing at 2.3063544244 yr, stable \(\to\) unstable. \\
\bottomrule
\end{tabularx}
\end{table}

At the exponential-update mobilising crossing, the unit-modulus pair is approximately \(0.9846404664\pm0.1745942495i\), with angle 0.175494 radians. The exponential-update crossing estimate is unchanged to the printed digits on 2,001-, 20,001-, and 200,001-point preliminary scans; each reported value is then refined by 70 bisection iterations. The sign of \(\rho(1)-1\) is only about \(3.5066\times10^{-4}\) at baseline. Euler crossings at 47.536 and 79.143 years, and the Euler protective \(-1\) crossing at 2.306 years, are properties of the Euler update; they are not numerical approximations to the exponential-update boundary at those intervals.

A continuous-delay comparator from the archived companion calculation has approximate crossings at 3.666 and 150.358 years. It is retained as an operator contrast, not as a third discretization of the same review map. The exponential frozen-assessment result and the continuous-delay result have different state/operator constructions; their separation cannot be attributed to a single parameter while changing all architecture at once.

\subsection{\texorpdfstring{Re-run one-at-a-time \(\pm2\%\) sensitivity}{Re-run one-at-a-time +/-2 percent sensitivity}}

\begin{table}[htbp]
\centering\scriptsize
\caption{Re-run first exponential-update mobilising crossing (years) after one-at-a-time changes in each parameter. ``None'' means no crossing was found on the registered \([0.2,200]\)-year scan.}
\label{tab:ota-sensitivity}
\begin{tabular}{@{}lrrrr@{}}
\toprule Parameter & \(-2\%\) & \(-1\%\) & \(+1\%\) & \(+2\%\) \\
\midrule
\(K\) & None & 2.986 & 8.665 & 10.371 \\
\(q\) & None & 3.438 & 8.502 & 10.100 \\
\(E_{\max}\) & None & 3.445 & 8.499 & 10.096 \\
\(\eta\) & 5.270 & 5.924 & 7.021 & 7.494 \\
\(r\) & 7.349 & 6.936 & 6.041 & 5.549 \\
\(\Delta_{\rm ref}\) & 10.775 & 8.891 & 2.353 & None \\
\(\tau_m\) & 6.599 & 6.551 & 6.449 & 6.394 \\
\(\delta_0\) & 6.495 & 6.498 & 6.505 & 6.508 \\
\(Z_{\rm ref}\) & 6.506 & 6.504 & 6.499 & 6.497 \\
\bottomrule
\end{tabular}
\end{table}

The most sensitive directions among this scalar battery are \(K,q,E_{\max},\Delta_{\rm ref}\). A missing crossing is a qualitative change in this scanned result, not simply a displaced point estimate. The signs are not symmetric: for example, a 2\% increase in \(\Delta_{\rm ref}\) removes the crossing from the scan, while a 2\% decrease moves it to 10.775 years. The table is deterministic parameter perturbation, not statistical uncertainty.

\subsection{Archived one-at-a-time \texorpdfstring{\(\pm5\%\)}{+/-5 percent} battery}

A separate archived sensitivity battery varied ten parameters at factors 0.90, 0.95, 1.05, and 1.10. It is retained as an archived record, not described as the current v53 rerun. The crossing-only summary is reproduced from \path{src/p7_empirical/recovered/sens_table.csv}; the file also retains equilibrium values, annual radii, crossing angles, crossing types, and protective-grid maxima. ``None'' denotes no crossing on the archived interval scan. The re-run \(\pm2\%\) table above and this archived wider perturbation record are distinct analyses.

\begin{table}[htbp]
\centering\small
\caption{Archived one-at-a-time \(\pm5\%\) and \(\pm10\%\) sensitivity battery. Entries are first exponential-update mobilising crossings in years; ``None'' means no crossing on the archived interval scan.}
\label{tab:ota5}
\begin{tabular}{@{}lrrrr@{}}
\toprule
Parameter & \(-10\%\) & \(-5\%\) & \(+5\%\) & \(+10\%\) \\
\midrule
\(r\) & 10.1844 & 8.4928 & 3.7490 & None \\
\(K\) & None & None & 14.2825 & 19.0977 \\
\(q\) & None & None & 13.7973 & 18.3769 \\
\(E_{\max}\) & None & None & 13.7899 & 18.3665 \\
\(\eta\) & None & 2.1944 & 8.7154 & 10.3140 \\
\(\delta_0\) & 6.4676 & 6.4845 & 6.5181 & 6.5348 \\
\(Z_{\rm ref}\) & 6.5268 & 6.5136 & 6.4897 & 6.4788 \\
\(\Delta_{\rm ref}\) & 20.8565 & 15.1725 & None & None \\
\(\tau_m\) & 6.9012 & 6.7279 & 6.2138 & 5.8545 \\
\(\delta\) & None & 4.6883 & 7.6905 & 8.5444 \\
\bottomrule
\end{tabular}
\end{table}


\subsection{Six-parameter joint perturbation: 20,001-point rerun}

The six jointly perturbed parameters are \((r,K,q,\eta,E_{\max},\Delta_{\rm ref})\), each independently set to the lower or upper edge of its stated relative box, producing 64 corners per box. Both boxes were recomputed from scratch on 20,001 evenly spaced review intervals in \([0.2,200]\) years (grid spacing about 0.00999 years) and crossings were bisected. The full per-corner JSON is retained as \texttt{\detokenize{work/joint_box_fullgrid.json}}.

\begin{table}[htbp]
\centering\small
\caption{Full-grid joint-box results. ``Stable no crossing'' means \(\rho<1\) at all 20,001 points; it is not a proof of stability on the continuum.}
\label{tab:joint-box}
\begin{tabular}{@{}lrrrrr@{}}
\toprule Box & Corners & Crossing & No crossing & Crossing range (yr) & Max \(\rho\) among no-crossing corners \\
\midrule
\(\pm0.5\%\) & 64 & 60 & 4 & 0.871965--10.670425 & 0.9999712801 \\
\(\pm1\%\) & 64 & 44 & 20 & 4.423192--13.653242 & 0.9999952956 \\
\bottomrule
\end{tabular}
\end{table}

Every no-crossing corner in both boxes is stable at all evaluated grid points; no no-crossing corner is grid-unstable. Because the maximum radii are close to one, the status is numerically delicate. These boxes are not posterior or confidence intervals, and no sampling distribution is assigned to the parameters. The results show that a single baseline crossing should not be reported as a robust institutional threshold.

\subsection{Independent numerical and nonlinear checks}

The retained precision-certificate log reports a high-precision root \(T_{\mathrm r,c}=6.50129494807546836889974861684\), consistent with the double-precision monodromy crossing. A fresh independent audit of the full nonlinear plant and nonlinear frozen-assessment controller flow is in \texttt{\detokenize{work/nonlinear_flow_check_v53.py}} and its JSON output. With DOP853 integration and central differences, the full-map Jacobian agrees with the monodromy to maximum absolute error \(1.67\times10^{-9}\) at the crossing for relative step \(3\times10^{-5}\); finite-difference crossing estimates across five steps range from 6.501276 to 6.501299 years. The prior normal-form script used a linearized controller update and therefore does not establish a normal-form coefficient for the full nonlinear map; its supercritical label is not relied on here. No invariant-circle continuation or independent orbit confirmation was completed, so no nonlinear bifurcation theorem is claimed.

\subsection{F-grid placement convention}

The model's equilibrium fishing-mortality proxy is swept as \(F^*=qE^*\) within its illustrative parameter grid. For the cross-system placement summaries in S13.2, the code uses fixed categories: below-stable if \(F<0.001\); marginal if \(0.001\leq F\leq0.004\); unstable-range if \(0.004<F\leq0.02\); beyond-feasible if \(F>0.02\). These names describe placement relative to the model grid and its maximum intrinsic growth rate, not assessed stock status. Their ecological scale is not calibrated to the 42-stock cohort; treat the categories as exploratory diagnostics only.

